% GENERATED FILE. Edit canonical lessons, then run npm run generate:books.
% canonical-source-hash: fnv1a64:f3bfafd15e4baf05
% canonical-lessons: ML-C224-kalarttuka, ML-C224-vilambuka, ML-C224-rucikkuka, ML-C224-niraykkuka, ML-C224-mutuka

\chapter{To mix, To serve, To taste, To fill, To cover}
\label{ch:ml-to-mix-to-serve-to-taste-to-fill-to-cover}
\hlchaptermodality{224}

\begin{chapteropening}
\textbf{By the end of this chapter:} \emph{I can say to mix, to serve, to taste, to fill and to cover.}

Complete the last lesson of chapter 224: I can say to mix, to serve, to taste, to fill and to cover.
\end{chapteropening}

\section[\texorpdfstring{\ml{കലർത്തുക} (kalarttuka)}{കലർത്തുക (kalarttuka)}]{\ml{കലർത്തുക} (kalarttuka) — to mix}

\label{lesson:ML-C224-kalarttuka}



\begin{quote}
Before the new one: say the Malayalam for to take off, then the Malayalam for to change.
\end{quote}

\subsection*{You'll want to know: \ml{കലർത്തുക}}
\textbf{\ml{കലർത്തുക}} — \emph{kalarttuka} — \textquotedblleft{}to mix\textquotedblright{}.

\begin{grammarlens}[title={What you've built}]
One more word for this chapter.
\end{grammarlens}

\subsection*{Your turn}
\begin{itemize}
\raggedright
  \item \emph{Say it:} \emph{kalarttuka}
  \item \emph{Say it:} \emph{kalarttuka}, once more
  \item \emph{Say it:} say \emph{kalarttuka}
  \item \emph{Recall:} say the Malayalam for to take off, then the Malayalam for to change, then say \emph{kalarttuka} again
\end{itemize}

\begin{tcolorbox}[breakable,colback=teal!4,colframe=teal!35!black,title={Before you move on}]
What does \ml{കലർത്തുക} mean? (\textquotedblleft{}To mix\textquotedblright{}.) Say it once more.
\end{tcolorbox}
\section[\texorpdfstring{\ml{വിളമ്പുക} (viḷambuka)}{വിളമ്പുക (viḷambuka)}]{\ml{വിളമ്പുക} (viḷambuka) — to serve (food)}

\label{lesson:ML-C224-vilambuka}



\begin{quote}
Before the new one: say the Malayalam for to change, then the Malayalam for to mix.
\end{quote}

\subsection*{You'll want to know: \ml{വിളമ്പുക}}
\textbf{\ml{വിളമ്പുക}} — \emph{viḷambuka} — \textquotedblleft{}to serve (food)\textquotedblright{}.

\begin{grammarlens}[title={What you've built}]
One more word for this chapter.
\end{grammarlens}

\subsection*{Your turn}
\begin{itemize}
\raggedright
  \item \emph{Say it:} \emph{viḷambuka}
  \item \emph{Say it:} \emph{viḷambuka}, once more
  \item \emph{Say it:} say \emph{viḷambuka}
  \item \emph{Recall:} say the Malayalam for to change, then the Malayalam for to mix, then say \emph{viḷambuka} again
\end{itemize}

\begin{tcolorbox}[breakable,colback=teal!4,colframe=teal!35!black,title={Before you move on}]
What does \ml{വിളമ്പുക} mean? (\textquotedblleft{}To serve (food)\textquotedblright{}.) Say it once more.
\end{tcolorbox}
\section[\texorpdfstring{\ml{രുചിക്കുക} (rucikkuka)}{രുചിക്കുക (rucikkuka)}]{\ml{രുചിക്കുക} (rucikkuka) — to taste}

\label{lesson:ML-C224-rucikkuka}



\begin{quote}
Before the new one: say the Malayalam for to mix, then the Malayalam for to serve.
\end{quote}

\subsection*{You'll want to know: \ml{രുചിക്കുക}}
\textbf{\ml{രുചിക്കുക}} — \emph{rucikkuka} — \textquotedblleft{}to taste\textquotedblright{}.

\begin{grammarlens}[title={What you've built}]
One more word for this chapter.
\end{grammarlens}

\subsection*{Your turn}
\begin{itemize}
\raggedright
  \item \emph{Say it:} \emph{rucikkuka}
  \item \emph{Say it:} \emph{rucikkuka}, once more
  \item \emph{Say it:} say \emph{rucikkuka}
  \item \emph{Recall:} say the Malayalam for to mix, then the Malayalam for to serve, then say \emph{rucikkuka} again
\end{itemize}

\begin{tcolorbox}[breakable,colback=teal!4,colframe=teal!35!black,title={Before you move on}]
What does \ml{രുചിക്കുക} mean? (\textquotedblleft{}To taste\textquotedblright{}.) Say it once more.
\end{tcolorbox}
\section[\texorpdfstring{\ml{നിറയ്ക്കുക} (niṟaykkuka)}{നിറയ്ക്കുക (niṟaykkuka)}]{\ml{നിറയ്ക്കുക} (niṟaykkuka) — to fill}

\label{lesson:ML-C224-niraykkuka}



\begin{quote}
Before the new one: say the Malayalam for to serve, then the Malayalam for to taste.
\end{quote}

\subsection*{You'll want to know: \ml{നിറയ്ക്കുക}}
\textbf{\ml{നിറയ്ക്കുക}} — \emph{niṟaykkuka} — \textquotedblleft{}to fill\textquotedblright{}.

\begin{grammarlens}[title={What you've built}]
One more word for this chapter.
\end{grammarlens}

\subsection*{Your turn}
\begin{itemize}
\raggedright
  \item \emph{Say it:} \emph{niṟaykkuka}
  \item \emph{Say it:} \emph{niṟaykkuka}, once more
  \item \emph{Say it:} say \emph{niṟaykkuka}
  \item \emph{Recall:} say the Malayalam for to serve, then the Malayalam for to taste, then say \emph{niṟaykkuka} again
\end{itemize}

\begin{tcolorbox}[breakable,colback=teal!4,colframe=teal!35!black,title={Before you move on}]
What does \ml{നിറയ്ക്കുക} mean? (\textquotedblleft{}To fill\textquotedblright{}.) Say it once more.
\end{tcolorbox}
\section[\texorpdfstring{\ml{മൂടുക} (mūṭuka)}{മൂടുക (mūṭuka)}]{\ml{മൂടുക} (mūṭuka) — to cover}

\label{lesson:ML-C224-mutuka}



\begin{quote}
Before the new one: say the Malayalam for to taste, then the Malayalam for to fill.
\end{quote}

\subsection*{You'll want to know: \ml{മൂടുക}}
\textbf{\ml{മൂടുക}} — \emph{mūṭuka} — \textquotedblleft{}to cover\textquotedblright{}.

\begin{grammarlens}[title={What you've built}]
One more word for this chapter.
\end{grammarlens}

\subsection*{Your turn}
\begin{itemize}
\raggedright
  \item \emph{Say it:} \emph{mūṭuka}
  \item \emph{Say it:} \emph{mūṭuka}, once more
  \item \emph{Say it:} say \emph{mūṭuka}
  \item \emph{Recall:} say the Malayalam for to taste, then the Malayalam for to fill, then say \emph{mūṭuka} again
\end{itemize}

\begin{tcolorbox}[breakable,colback=teal!4,colframe=teal!35!black,title={Before you move on}]
What does \ml{മൂടുക} mean? (\textquotedblleft{}To cover\textquotedblright{}.) Say it once more.
\end{tcolorbox}
